\section{DSM-III was never clustered into existence}\label{11_dsm_iii-dsm-iii-was-never-clustered-into-existence}

\textbf{DSM-III's categories were not derived by factor analysis, cluster analysis, or any multivariate method --- they were produced by committee consensus and voting, with kappa statistics applied only after the fact to measure interrater agreement.} This single historical fact reframes every serious argument against the manual: the ``statistical methodology'' of DSM-III is largely a myth. The authoritative nosology of modern psychiatry was built on the neo-Kraepelinian intuitions of roughly a dozen clinicians at Washington University and New York State Psychiatric Institute, formalized into polythetic criteria, and then defended with a rhetorical move --- \emph{reliability is a necessary condition for validity} --- that was never followed up by any actual validity test. The manifesto's argument that clustering methods were inappropriate for psychiatric symptoms is therefore even stronger than it appears: clustering methods were not merely misapplied, they were \textbf{assumed in the format of the criteria without ever being run}. Every subsequent empirical challenge --- comorbidity rates of 45--66\%, 636,120 ways to qualify for PTSD, \href{https://www.researchgate.net/publication/260750182_636120_Ways_to_Have_Posttraumatic_Stress_Disorder}{ResearchGate} a single \emph{p}-factor better-fitting than eleven distinct disorders, cross-disorder genetic correlations of 0.43--0.68 --- is best understood not as falsifying a carefully derived taxonomy but as exposing the epistemic emptiness of a taxonomy that was never derived at all.

\subsection{The statistical method that wasn't}\label{11_dsm_iii-the-statistical-method-that-wasnt}

The historical record, most thoroughly documented in Hannah Decker's \emph{The Making of DSM-III} (Oxford, 2013) and in Davies's 2016 \emph{Anthropology and Medicine} interviews, is unambiguous. Spitzer himself described the process: ``\textbf{We simply discussed things until we were comfortable with it\ldots{} It was done by consensus of experts, which would now be considered a very trivial approach, but that's the best we could do then.}'' \href{https://www.tandfonline.com/doi/full/10.1080/13648470.2016.1226684}{Taylor \& Francis Online} The Feighner criteria (1972) --- the direct ancestors of DSM-III --- were drafted by a Washington University discussion group \href{https://pubmed.ncbi.nlm.nih.gov/20008944/}{PubMed} starting from prior clinical papers like Cassidy et al.'s depression symptom list and modifying them through conversation. The famous 6-of-10 threshold for depression was chosen, in group member Muñoz's words, ``\textbf{because it sounded about right.}'' Kendler, Aggen and Neale (\emph{JAMA Psychiatry} 2013) confirm: ``Clinical judgment and not psychometrics guided the development of these criteria.'' \href{https://pmc.ncbi.nlm.nih.gov/articles/PMC3800168/}{PubMed Central}

The only statistic that played a central role was \textbf{Cohen's/Fleiss's kappa}, operationalized in Spitzer \& Fleiss's 1974 \emph{British Journal of Psychiatry} re-analysis. That paper's rhetorical move is the pivot of the entire DSM enterprise: ``\emph{A necessary constraint on the validity of a system is its reliability. There is no guarantee that a reliable system is valid, but assuredly an unreliable system must be invalid.}'' \href{https://www.cambridge.org/core/journals/the-british-journal-of-psychiatry/article/abs/reanalysis-of-the-reliability-of-psychiatric-diagnosis/FEBECBB1AEB68DA969F897333B8003E3}{Cambridge Core} Having defined pre-1974 diagnosis as unacceptably unreliable, Spitzer created the crisis DSM-III would solve. \href{https://academyanalyticarts.org/kirk-myth-reliability-dsm}{Academyanalyticarts} The DSM-III field trials (Spitzer, Forman, Nee, \emph{AJP} 1979) then reported Axis I kappa of 0.78 --- but the design was methodologically favorable: joint interviews (where two raters share the same observation stream inflate agreement), volunteer sympathetic clinicians, and aggregated statistics hiding individual-diagnosis kappas below 0.5. \href{https://pubmed.ncbi.nlm.nih.gov/443467/}{PubMed} Chmielewski, Clark, Bagby \& Watson (\emph{Journal of Abnormal Psychology} 2015) showed that under realistic test-retest independent interviews, mean DSM-IV kappa was \textasciitilde0.47 --- fair, not good. \href{https://pmc.ncbi.nlm.nih.gov/articles/PMC4573819/}{PubMed Central} Kraemer et al.~(2012) quietly relaxed what would count as ``acceptable'' kappa for DSM-5 rather than confront these numbers. \href{https://www.researchgate.net/publication/264796782_Reliability_in_Psychiatric_Diagnosis_with_the_DSM_Old_Wine_in_New_Barrels}{ResearchGate}

The embedded assumptions were never tested, only assumed: that disorders are \textbf{discrete categorical kinds} (no taxometric check); that symptoms are \textbf{exchangeable compensatory indicators} of those kinds (the 5-of-9 polythetic rule assumes equal weighting); that \textbf{interrater reliability is a precondition for validity} rather than a measure of rule-application agreement; and that \textbf{committee consensus is epistemically equivalent to empirical discovery}. The Robins-Guze (1970) validator framework --- clinical description, laboratory study, delimitation, follow-up, family study --- was offered as the validation program, but phases 2-5 were never completed for most disorders.

\subsection{Three converging formal critiques}\label{11_dsm_iii-three-converging-formal-critiques}

Three lines of formal argument from recognized authorities independently converge on rejecting DSM validity.

\textbf{Kendler's mechanistic property cluster (MPC) framework} (Kendler, Zachar \& Craver, \emph{Psychological Medicine} 2011, 41:1143--1150) systematically rules out the four candidate ontologies of psychiatric disorder: essentialist natural kinds \href{https://scholar.google.com/scholar_lookup?title=What+kinds+of+things+are+psychiatric+disorders?&amp=&author=K.+S.+Kendler&amp=&author=P.+Zachar&amp=&author=C.+Craver&amp=&publication_year=2011&amp=&pages=1143-1150&amp=&doi=10.1017/S0033291710001844}{Google Scholar} (falsified by polygenicity and heterogeneity), social constructs (falsified by cross-cultural recurrence), \textbf{practical kinds} (the DSM's implicit stance, rejected as scientifically unambitious because it ``does not require that diagnoses be grounded in shared causal processes''), and MPC kinds (endorsed). \href{https://scholar.google.com/scholar_lookup?title=What+kinds+of+things+are+psychiatric+disorders?&amp=&author=K.+S.+Kendler&amp=&author=P.+Zachar&amp=&author=C.+Craver&amp=&publication_year=2011&amp=&pages=1143-1150&amp=&doi=10.1017/S0033291710001844}{Google Scholar} Under MPC, disorders are ``fuzzy sets defined by mechanisms at multiple levels\ldots{} populations with central paradigmatic and more marginal members.'' \href{https://scholar.google.com/scholar_lookup?title=What+kinds+of+things+are+psychiatric+disorders?&amp=&author=K.+S.+Kendler&amp=&author=P.+Zachar&amp=&author=C.+Craver&amp=&publication_year=2011&amp=&pages=1143-1150&amp=&doi=10.1017/S0033291710001844}{Google Scholar} Because mechanisms --- not surface features --- define the kinds, \textbf{symptom-level clustering cannot in principle recover valid categories}. Kendler's 2012 \emph{Molecular Psychiatry} paper extends this with ``dappled'' causation: causes are distributed in irregular patches across biological, psychological, and social levels, so any flat symptom-level procedure is mechanism-blind.

\textbf{Hyman's reification critique} (\emph{Annual Review of Clinical Psychology} 2010, 6:155--179) introduces the most quotable formulation: DSM created ``an \textbf{unintended epistemic prison that was palpably impeding scientific progress}.'' \href{https://philarchive.org/archive/HYMTDO-3}{philarchive} Hyman documents that researchers proceeded ``as if {[}DSM disorders{]} were \textbf{natural kinds that would map onto the human genome}.'' \href{https://philarchive.org/archive/HYMTDO-3}{philarchive} They do not. Two patients with identical MDD diagnoses can share only one of nine criteria. \href{https://philarchive.org/archive/HYMTDO-3}{philarchive} His structural argument has five steps: reliability ≠ validity; polythetic criteria breed within-category heterogeneity; \href{https://www.annualreviews.org/doi/pdf/10.1146/annurev.clinpsy.3.022806.091532}{Annual Reviews} comorbidity of \textasciitilde50\% suggests miscarved boundaries; family and genetic data disrespect DSM categories; and FDA/grant/journal lock-in makes the reification self-sustaining.

\textbf{Insel's NIMH announcement} (Director's Blog, 29 April 2013) was the institutional rupture. ``DSM has been described as a `Bible' for the field, \textbf{it is, at best, a dictionary}\ldots{} \textbf{The weakness is its lack of validity.} Unlike our definitions of ischemic heart disease, lymphoma, or AIDS, DSM diagnoses are based on a consensus about clusters of clinical symptoms, not any objective laboratory measure.'' \href{https://psychrights.org/2013/130429NIMHTransformingDiagnosis.htm}{Psychrights} Crucially: ``\textbf{We cannot succeed if we use DSM categories as the `gold standard'\ldots{} Imagine deciding that EKGs were not useful because many patients with chest pain did not have EKG changes.}'' \href{https://psychrights.org/2013/130429NIMHTransformingDiagnosis.htm}{Psychrights} The RDoC framework (Insel et al., \emph{AJP} 2010; Cuthbert \& Insel, \emph{BMC Medicine} 2013) operationalizes this into dimensional, mechanism-based units of analysis. \href{https://link.springer.com/article/10.1186/1741-7015-11-126}{Springer} The convergence with Kendler and Hyman is explicit: RDoC cites both.

\subsection{The empirical wreckage}\label{11_dsm_iii-the-empirical-wreckage}

The numerical evidence against DSM categorical validity has reached the point of absurdity.

\textbf{Comorbidity runs on a rule of 50\%}. The National Comorbidity Survey Replication (Kessler et al., \emph{Archives of General Psychiatry} 2005) \href{https://pmc.ncbi.nlm.nih.gov/articles/PMC2847357/}{PubMed Central} found 45\% of people with any 12-month DSM-IV disorder met criteria for two or more. The Dunedin birth-cohort data analyzed by Caspi \& Moffitt (\emph{AJP} 2018) show the pattern is fractal: of those meeting criteria for one disorder, 66\% met criteria for a second; of those with two, 53\% met criteria for a third; of those with three, 41\% met criteria for a fourth. \href{https://pubmed.ncbi.nlm.nih.gov/29621902/}{PubMed} \textbf{Diagnostic stability collapses over time}: Bromet et al.~(\emph{AJP} 2011) followed first-admission psychosis for ten years and found \textbf{50.7\% of patients changed diagnostic category}; \href{https://pmc.ncbi.nlm.nih.gov/articles/PMC3589618/}{PubMed Central} Suárez-Pinilla et al.~reported only 47.7\% stability at 21 years in the PAFIP cohort, with 80\% having changed by year 10. \href{https://www.cambridge.org/core/journals/psychological-medicine/article/abs/longterm-diagnostic-stability-predictors-of-diagnostic-change-and-time-until-diagnostic-change-of-firstepisode-psychosis-a-21year-followup-study/2D7A6A5775A2AEA714260DC85BB4A19D}{Cambridge Core}

\textbf{Within-category heterogeneity explodes combinatorially}. Fried \& Nesse (\emph{Journal of Affective Disorders} 2015) analyzed STAR\emph{D (n=3,703) \href{https://www.sciencedirect.com/science/article/abs/pii/S0165032714006326}{ScienceDirect} and found \textbf{1,030 unique symptom profiles} for DSM-defined MDD; \href{https://www.sciencedirect.com/science/article/abs/pii/S027273582100043X}{ScienceDirect} 83.9\% of profiles were endorsed by five or fewer participants. \href{https://asu.elsevierpure.com/en/publications/depression-is-not-a-consistent-syndrome-an-investigation-of-uniqu/}{Arizona State University} With decompounded sleep/appetite/psychomotor items, MDD admits \textbf{10,377 qualifying combinations}, \href{https://pmc.ncbi.nlm.nih.gov/articles/PMC4397113/}{PubMed Central} \href{https://www.sciencedirect.com/science/article/pii/S0278584620304498}{ScienceDirect} rising to \textbf{341,737 with the melancholic specifier}. \href{https://eiko-fried.com/10377-ways-for-major-depression-but-341737-ways-for-melancholia/}{Eiko Fried +2} Galatzer-Levy \& Bryant's }Perspectives on Psychological Science* (2013) paper --- titled ``\textbf{636,120 Ways to Have Posttraumatic Stress Disorder}'' --- documents that DSM-5 expanded PTSD's polythetic combinatorics eightfold from DSM-IV. \href{https://journals.sagepub.com/doi/abs/10.1177/1745691613504115}{Sage Journals +4} Olbert, Gala \& Tupler (\emph{Journal of Abnormal Psychology} 2014) proved formally that in most DSM categories, \textbf{two patients with the same diagnosis can share zero symptoms}. \href{https://pubmed.ncbi.nlm.nih.gov/24886017/}{PubMed +2}

\textbf{Treatment response betrays categorical specificity in both directions}. Within MDD, STAR\emph{D's (}AJP* 2006) sequential remission rates were 36.8\%, 30.6\%, 13.7\%, 13.0\% across steps; \href{https://link.springer.com/article/10.1007/s11920-007-0061-3}{Springer} Pigott et al.'s 2023 \emph{BMJ} reanalysis found the true cumulative remission was \textasciitilde35\%, not the advertised 67\%. \href{https://www.psychiatrictimes.com/view/star-d-s-cumulative-remission-rate-and-why-it-still-matters}{Psychiatry Times} No active antidepressant beat any other. \href{https://journals.sagepub.com/doi/pdf/10.1177/070674371005500303}{Sage Journals} CATIE (Lieberman et al., \emph{NEJM} 2005) found 74\% of schizophrenia patients discontinued their assigned antipsychotic within 18 months, \href{https://pmc.ncbi.nlm.nih.gov/articles/PMC6482974/}{PubMed Central} \href{https://pmc.ncbi.nlm.nih.gov/articles/PMC5935444/}{PubMed Central} with no advantage for newer atypicals. \href{https://psychiatryonline.org/doi/10.1176/appi.ajp.2011.11010039}{Psychiatry Online} Across categories, \textbf{the same drugs work for ostensibly distinct disorders}: SSRIs are first-line for MDD, GAD, panic, OCD, PTSD, social anxiety, and PMDD; lithium reduces suicidality across diagnostic boundaries; Barlow's Unified Protocol (\emph{JAMA Psychiatry} 2017) matched diagnosis-specific CBT across GAD, SAD, OCD, and panic. \href{https://onlinelibrary.wiley.com/doi/10.1002/wps.20748}{Wiley Online Library} Cipriani et al.'s 522-trial network meta-analysis (\emph{Lancet} 2018) found all 21 antidepressants modestly superior to placebo (SMD 0.30) with no meaningful specificity. \href{https://pubmed.ncbi.nlm.nih.gov/29477251/}{PubMed} \href{https://research-information.bris.ac.uk/en/publications/comparative-efficacy-and-acceptability-of-21-antidepressant-drugs/}{University of Bristol}

\textbf{Genetics refuses to respect DSM boundaries}. The Cross-Disorder Psychiatric Genomics Consortium (Lee et al., \emph{Nature Genetics} 2013) estimated SNP-based genetic correlations \href{https://www.thelancet.com/journals/lancet/article/PIIS0140-6736(12)62129-1/abstract}{The Lancet} \href{https://scholar.google.com/scholar_lookup?hl=en&volume=381&publication_year=2013&pages=1371-1379&journal=Lancet&author=Cross-Disorder+Group+of+the+Psychiatric+Genomics+Consortium&title=Identification+of+risk+loci+with+shared+effects+on+five+major+psychiatric+disorders:+A+genome-wide+analysis}{Google Scholar} of \textbf{0.68 between schizophrenia and bipolar disorder, 0.47 bipolar-MDD, 0.43 schizophrenia-MDD}. The Brainstorm Consortium (\emph{Science} 2018), analyzing 265,218 patients across 25 brain disorders, found psychiatric disorders share common-variant risk promiscuously while neurological disorders are genetically distinct --- \textbf{psychiatric DSM categories do not cut at genetic joints}. \href{https://kclpure.kcl.ac.uk/portal/en/publications/analysis-of-shared-heritability-in-common-disorders-of-the-brain/}{King's College London} \href{https://cea.hal.science/cea-01870483v1/document}{Hal} Selzam et al.~(\emph{Translational Psychiatry} 2018) extracted a \textbf{polygenic p-factor explaining 22-57\% of variance} across methods, \href{https://www.nature.com/articles/s41398-018-0217-4}{Nature} paralleling Caspi et al.'s (\emph{Clinical Psychological Science} 2014) behavioral p-factor that fit Dunedin data better than three-factor or category models. \href{https://pubmed.ncbi.nlm.nih.gov/25360393/}{PubMed} \href{https://pmc.ncbi.nlm.nih.gov/articles/PMC4209412/}{PubMed Central} Kendler's own 1996 twin work established that GAD and MDD share essentially identical genetic liabilities.

\subsection{The dimensional alternative and its mathematical teeth}\label{11_dsm_iii-the-dimensional-alternative-and-its-mathematical-teeth}

The Hierarchical Taxonomy of Psychopathology (Kotov et al., \emph{Journal of Abnormal Psychology} 2017, 126:454--477) is the consensus document of 40 leading nosologists specifying what a post-DSM taxonomy looks like: \href{https://pubmed.ncbi.nlm.nih.gov/28333488/}{PubMed} \href{https://class.unt.edu/hitop/images/hitop.unt.edu/files/kotov_et_al_2017_-_hitop.pdf}{Unt} six empirical spectra (internalizing, thought disorder, disinhibited externalizing, antagonistic externalizing, detachment, somatoform) organized into three superspectra under a general \emph{p}-factor, with syndromes below and symptom components at the bottom. \href{https://pmc.ncbi.nlm.nih.gov/articles/PMC6859953/}{PubMed Central} Its methods --- confirmatory factor analysis, bifactor modeling, Goldberg's bass-ackwards procedure, latent-class-vs-latent-trait comparisons across samples up to 43,093 participants --- represent what DSM-III could have done in 1980 and did not. \href{https://www.apa.org/pubs/journals/features/abn-abn0000258.pdf}{APA} \href{https://www.researchgate.net/publication/315591130_The_Hierarchical_Taxonomy_of_Psychopathology_HiTOP_A_Dimensional_Alternative_to_Traditional_Nosologies}{ResearchGate} Watson et al.~(\emph{World Psychiatry} 2022) state the verdict plainly: ``\textbf{not a single mental disorder has been established in the scientific literature as a discrete categorical entity}.'' \href{https://en.wikipedia.org/wiki/Hierarchical_Taxonomy_of_Psychopathology}{Wikipedia}

The formal mathematical case against DSM's categorical format runs through \textbf{MacCallum, Zhang, Preacher \& Rucker} (\emph{Psychological Methods} 2002), whose worked examples show \href{https://www.semanticscholar.org/paper/On-the-practice-of-dichotomization-of-quantitative-Maccallum-Zhang/67646d3def337c40bbfa2476f0f1ff785ae18242}{Semantic Scholar} dichotomization of a continuous variable with true correlation r = .30 reduces the observed correlation to \textasciitilde.21, discarding more than half the explained variance. \href{https://quantpsy.org/pubs/maccallum_zhang_preacher_rucker_2002.pdf}{Quantpsy} Applied to DSM's 5-of-9 MDD threshold or 6-of-20 PTSD threshold, this is a rigorous demonstration that the polythetic cutoff \textbf{destroys information rather than compressing it usefully}. Haslam, McGrath, Viechtbauer \& Kuppens's (\emph{Psychological Medicine} 2020) meta-analysis of 317 taxometric findings using the Comparative Curve Fit Index found dimensional models outperforming taxonic models \textbf{5:1}, with zero content domains aggregating as taxonic \href{https://cris.maastrichtuniversity.nl/en/publications/dimensions-over-categories-a-meta-analysis-of-taxometric-research/}{Maastricht University} --- only six variables (alcohol dependence, IED, problem gambling, autism, suicide risk, pedophilia) showed replicated taxonic signatures. \href{https://pubmed.ncbi.nlm.nih.gov/32493520/}{PubMed}

Borsboom's network theory (\emph{World Psychiatry} 2017, 16:5--13) supplies the alternative ontology: \href{https://dennyborsboom.com/publications/}{Denny Borsboom} \textbf{disorders are not latent entities causing symptoms but alternative stable states of causally connected symptom networks}. \href{https://pubmed.ncbi.nlm.nih.gov/28127906/}{PubMed} When connectivity exceeds a threshold, symptom activation becomes self-sustaining. \href{https://pubmed.ncbi.nlm.nih.gov/23537483/}{PubMed +2} This dispenses with the latent-variable assumption underlying both categorical and dimensional models that treat symptoms as indicators of some hidden construct.

\subsection{Brown clustering: a defensible analogy, but not the load-bearing one}\label{11_dsm_iii-brown-clustering-a-defensible-analogy-but-not-the-load-bearing-one}

No prior psychiatric paper uses Brown clustering as an analogy for DSM failure --- searches across PubMed, PMC, ACL Anthology, and computational psychiatry venues return nothing. The analogy is \textbf{defensible if handled carefully, cringe if handled casually}. Its strength is that Brown clustering's known failure mode is structurally identical to the DSM problem: it brilliantly groups \texttt{\{Monday,\ Tuesday,\ Wednesday,\ ...\}} (shallow distributional near-synonyms) but produces clusters like \texttt{\{great,\ big,\ vast,\ sudden,\ mere,\ sheer,\ gigantic,\ lifelong,\ scant,\ colossal\}} --- real in co-occurrence statistics, incoherent in semantic mechanism. \href{https://simonsuster.github.io/talks/Brown_RUG.pdf}{Simonsuster} That mirrors DSM: insomnia, fatigue, and low mood co-occur in clinic regardless of whether the generative mechanism is sleep apnea, grief, SSRI withdrawal, or hypothyroidism. A greedy agglomerative mutual-information algorithm \href{https://ojs.aaai.org/index.php/AAAI/article/download/10190/10049}{AAAI} with no theoretical guarantees \href{https://en.wikipedia.org/wiki/Brown_clustering}{Wikipedia} is exactly the right formal cousin of a committee-consensus symptom-grouping procedure.

The risks are that computational psychiatry audiences (Huys, Gillan, Stephan, Friston circles) may find Brown clustering arbitrary --- it's a 1992 algorithm long superseded by word2vec/GloVe/BERT --- and read its invocation as namedropping from someone who hasn't read the field's own critiques. The manifesto should \textbf{frame Brown clustering as a vivid illustrative cousin, not the primary frame}, and embed it inside the field's canonical vocabulary: Borsboom, Mellenbergh \& van Heerden (\emph{Psychological Review} 2004) on validity requiring ``(a) the attribute exists and (b) variations in the attribute causally produce variation in the measurement outcomes''; \href{https://dennyborsboom.com/wp-content/uploads/2017/01/borsboomvalidity2004.pdf}{Dennyborsboom} Michell's charge that psychometrics is ``pathological science'' for assuming measurability rather than demonstrating it (\emph{Theory \& Psychology} 2000; \emph{Measurement} 2008); and the computational-psychiatry consensus --- Huys, Maia \& Frank (\emph{Nature Neuroscience} 2016); \href{https://www.tnu.ethz.ch/fileadmin/user_upload/documents/Publications/2016/2016_Huys_Maia_Frank.pdf}{Ethz} \href{https://www.semanticscholar.org/paper/Computational-psychiatry-as-a-bridge-from-to-Huys-Maia/42abde7800c25b90ec1bb9c88dfbe2ab37e5ae16}{Semantic Scholar} Stephan et al.~(\emph{Lancet Psychiatry} 2016); Wise, Robinson \& Gillan (\emph{Biol Psychiatry} 2023) --- that the field has already abandoned DSM categories as research endpoints.

Stronger or complementary computational analogies include \textbf{phlogiston} (reliable descriptive compression, zero mechanism), \textbf{Ptolemaic epicycles} (DSM specifiers and NOS categories as structural epicycles), \textbf{k-means on non-convex manifolds} (the algorithm always returns clusters whether or not cluster structure exists), and \textbf{overfitting to noise} (Wiecki, Poland \& Frank, \emph{Clinical Psychological Science} 2015). The most intellectually robust framing combines them: DSM categories are to psychiatric mechanisms what Brown clusters are to semantic concepts, what phlogiston is to oxidation, and what epicycles are to gravity --- \textbf{reliable surface compressions of phenomena whose generative structure lies elsewhere}.

\subsection{Conclusion: the unclaimed information-theoretic argument}\label{11_dsm_iii-conclusion-the-unclaimed-information-theoretic-argument}

The strongest novel contribution a manifesto can make is not another rehearsal of the empirical wreckage but a \textbf{formal information-theoretic framing that the literature has not yet claimed}. Three gaps are open. First, \textbf{no one has formalized DSM categorization as a failed minimum description length (MDL) problem}: a valid categorical compression should produce shorter descriptions than the raw symptom vector while preserving predictive power over outcomes and treatment response. DSM labels manifestly fail both conditions --- 1,030 MDD profiles cannot be losslessly compressed to one symbol, \href{https://www.sciencedirect.com/science/article/abs/pii/S027273582100043X}{ScienceDirect +3} and the symbol does not predict SSRI response better than the vector. Second, \textbf{the ``encoding loss chain''} --- lived experience → patient's metaphor → clinician's interpretation → binary criterion checkbox --- is extensively described qualitatively (Kleinman's illness-vs-disease distinction; interview reliability literature) but has not been modeled as a cascade of information-destroying channels with quantifiable per-stage entropy loss. Third, \textbf{the disanalogy with linguistic corpora is sharp and unexploited}: Brown clustering works for language because \href{http://sean-chester.github.io/assets/preprints/generalised-brown.pdf}{Sean-chester} speakers share a generative grammar and stable distributional semantics; psychiatric symptoms lack this shared generative structure because the same surface symptom arises from mechanistically unrelated causal networks (multiple realizability, in Borsboom et al.'s 2019 \emph{BBS} formulation). The DSM assumed symptoms were like words in a corpus. They are not. The manifesto's deepest claim --- that clustering methods failed because psychiatric symptoms lack shared semantic structure --- is both correct and, as yet, unmade in the formal literature. That is the space to occupy.
